\chapter{The Rule System and Its Refusals}
\label{ch:refusals}

\section{Two Distinct Rule Sets}

A reader who goes looking for `the rules' will find two collections with
overlapping names, and confusing them makes the system incomprehensible. They
answer different questions.

The first is the \textbf{composer's law}: \LawRuleCount{} identifiers under
\LawFamilyCount{} headings, spanning \LawIdPrefixCount{} id prefixes, held in a
single document that is embedded in the application and parsed rather than
transcribed. It says what a well-formed CTW board \emph{is}: how wide a lane runs,
how far a wool room stands from a spawn, what a build zone may touch. It is
a design document with identifiers attached, and most of its rules are enforced
during composition, where the thing being judged still has a name.

The second is the \textbf{finding catalogue}: \RuleCount{} rules across
\RuleFamilyCount{} families, served live at \route{GET /api/rules}. It is the set of every rule that something,
somewhere, can raise as a complaint about a document. Its scope is the whole
studio rather than the composer, covering sketch geometry, dressing, house styles,
export and the request itself.

The two overlap without either containing the other. The composer's law carries
rules that nothing raises as a finding, and those are deliberately not served: the
question the catalogue exists to answer is \emph{what is this finding}, and a rule
nothing raises has no finding to explain. Measured against the running studio,
\LawServedCount{} of the law's \LawRuleCount{} ids are answered by the catalogue.
Everything else in the catalogue belongs to a gate the composer never touches.

\begin{table}[htbp]
\caption{The finding catalogue by family, generated from the running studio. A
family is the field the API answers with, not the id's prefix: the four
build-zone rules are numbered \id{BZ} and belong to family \id{CT}, so
\RuleIdPrefixCount{} prefixes span \RuleFamilyCount{} families.}
\label{tab:families-live}
\scriptsize
%% Generated from the running pgm-studio API (http://localhost:7894/api) --- do not edit by hand.
%% Source: GET /api/rules

\begin{longtable}{@{}L{16mm}R{12mm}L{112mm}@{}}
\toprule
\textbf{Family} & \textbf{Rules} & \textbf{What it governs} \\
\midrule
\endfirsthead
\toprule
\textbf{Family} & \textbf{Rules} & \textbf{What it governs} \\
\midrule
\endhead
\bottomrule
\endlastfoot
\id{SK} & 28 & The sketch document: shapes, layers, groups and mirroring, and the ground they rasterise into. \\
\id{DR} & 17 & Dressing a built world: where a tree, boulder, building or body of water may rest, and what it may not close off. \\
\id{PL} & 16 & The plan: pieces, placements, walls and the spawns and objectives a board is laid out around. \\
\id{WX} & 13 & The room frame stamped for a spawn or wool: shell, pad, door, iron, and where it stands on its piece. \\
\id{OB} & 12 & Objectives as they are played: where a goal may stand and float, what it is built of, and how the match can end. \\
\id{CT} & 10 & The middle: what the mid is carved into, the holes rotation leaves in it, and the build zones docked there (the BZ-numbered rules). \\
\id{HS} & 10 & House style: the materials and parts a building is authored from, part by part. \\
\id{WL} & 8 & Wool siting: a wool's distance from its spawn, its siblings and the frontline, and the approach into its room. \\
\id{EX} & 6 & The export gate: what must hold before a map document and its world are written. \\
\id{IM} & 6 & World import: the host fetched from, what arrived, and whether it holds a world the studio can read. \\
\id{RL} & 6 & Relief: whether the elevation is the landform the island claims, and whether it can be walked. \\
\id{RQ} & 6 & The request contract: which status a malformed, unknown, conflicting or unreadable request answers with. \\
\id{HJ} & 5 & Wing joints: the ways two rectangles of one building may meet, and the roof that follows. \\
\id{PT} & 5 & Terrain painting: how a theme states its materials, bands and sampled patterns. \\
\id{ST} & 5 & Structure pieces: what a wool-room or spawn piece defines, and the caps on what it may raise. \\
\id{FR} & 4 & The frontline: split or wide, how it docks, and how much of a face a crossing spans. \\
\id{GO} & 4 & Destroy-goal distances, measured by walk: to its own spawn, to the enemy's, and between goals. \\
\id{SP} & 4 & Spawn siting: where a spawn stands on its lane, and what its egress and its door face. \\
\id{DC} & 3 & Destroyables and cores: casing and interior, float against leak, and the material a size is worth. \\
\id{G} & 3 & The board envelope: corridor width, the void gaps a route hops, and map size against player count. \\
\id{HP} & 3 & The footprint a placed building is made of: how many rectangles, how thin, how much ground. \\
\id{ED} & 2 & The edit document: the references it must resolve, and the state it may be applied in. \\
\id{LN} & 2 & Lanes: the corridor width a board is built at, and the run before a junction or a dead end. \\
\id{CO} & 1 & Compose: descriptors that are well-formed and name a board the composer has no shape for. \\
\id{EL} & 1 & Land seams: the step a player cannot walk up, and the palette that steps by it. \\
\id{EZ} & 1 & Editable ground: standable ground nobody can edit, outside the zones a map seals on purpose. \\
\id{PC} & 1 & Piece connectivity: what counts as a join between two pieces, and what does not. \\
\id{SH} & 1 & Shops: the references a shop makes that the document has to define. \\
\midrule
\textbf{Total} & \textbf{183} & \textcolor{ink60}{28 families} \\
\end{longtable}

\end{table}

The full catalogue, rule by rule, is \cref{app:rules}.

\section{Anatomy of a Finding}

Every gate in the studio answers in one shape. A client rendering a refusal never
has to know which gate produced it before it can read one, and this uniformity is
the single most useful property the API has for a machine operator.

A finding is three things: a \textbf{rule id}, a \textbf{sentence}, and
\textbf{what it is about}.

The id is machine-legible and stable forever, so it outlives the task that added
it and never doubles as a task-tracking identifier. The sentence carries the
measured numbers, on an explicit principle. The word \emph{invalid} tells an
author nothing they can fix, so a refusal says \emph{this one is
2\,$\times$\,8} rather than \emph{this one is too small}. What it is about is a \textbf{field} where a
document field can be named, and \textbf{subjects} where the fault indicts things
an editor could highlight. A finding may carry either, both or neither.

One further field matters more to a machine than to a person. A finding states its
\textbf{edit} where the fix is mechanical: a document, a path, an operation and a
value. A reader applies it rather than re-deriving it from the rule's prose. It
is, as the documentation puts it, the half of a finding a model acts on where the
sentence alone is inspected and left.

What a finding deliberately does \emph{not} carry is the rule's classification.
The category and the concerns list are fixed by the id rather than by the raising
site, so they live on the rule and a caller joins them from the catalogue. Putting
them on the finding would put one fact at 96 sites.

\section{Three Severities}

This is the part of the system that has cost authoring runs the most, and it is
the first thing an agent should learn.

\begin{description}[leftmargin=1.6em,style=nextline,itemsep=4pt]
  \item[A refusal] stops the request. The work did not happen. It arrives in the
    error envelope's \id{findings}, and nothing rode along with it.
  \item[A decline] says that a piece of what was posted is \emph{not in what was
    built}. The request succeeded. Something you sent was dropped.
  \item[A complaint] is a remark. The request succeeded and the studio has an
    observation about it.
\end{description}

Both surviving severities ride on a success, under one key:

\begin{quote}
A 2xx JSON object answers \id{warnings} when something was complained about or
declined, and carries no such key when nothing was. \ldots A caller reads
\id{warnings ?? []} and is done.
\end{quote}

The consequence is blunt. A driver that reads only the status code throws away
half of every answer. A board can store at 200 with a floor missing under its
walls, with a theme that painted nothing, with a field it misspelled silently
unread, and the only place any of that is said is \id{warnings}. One run report
names this the dominant failure class in the whole corpus and calls it, exactly, a
200 that isn't a yes.

The distinction between the two is carried on each finding's \id{severity}, and
the documentation is explicit about why it could not be carried anywhere else:
\id{warnings.some(w => w.severity === "decline")} is \emph{something I sent was
dropped}, which no status code and no other field says.

\begin{ruling}{The Findings Read}
\route{GET /map/\{slug\}/findings} is asked on every run beside the others,
because it is the only read that answers \id{SK9}, the gate that knows two shapes
on one layer stacked and the lower one is not in the world. It is a
decline, the channel every other route publishes on keeps complaints alone, and so
a board can store at 200 with a floor missing under its walls and nothing anywhere
says so.
\end{ruling}

\section{The Refusal Envelope and Status Codes}

A refusal on the wire looks like this:

\begin{lstlisting}[language=json]
{ "error": "invalid house style",
  "message": "...",
  "findings": [ { "rule": "HS1", "message": "...",
                  "severity": "refusal", "field": "doorHead.block" } ] }
\end{lstlisting}

\id{error} names the \emph{gate}, such as \emph{objective placement},
\emph{unknown gamemode} or \emph{plan not compilable}, and never the fault itself.
\id{message} is the findings' sentences joined together.

Status codes stay the gate's own, and each means one thing: \textbf{400} a
document wrong as posted, \textbf{404} a subject the route names and the studio
does not have, \textbf{409} well-formed but conflicting with the map's state,
\textbf{422} cannot be processed, \textbf{500} the studio's own fault. Each route
declares which of them it can answer, 95 of 149 operations carrying at least one,
and a test fails if a documented row and its route stop agreeing.

The edge answers in the same shape as the gates. No route under \route{/api}
answers a failure in a shape of its own, and none answers one with no body at all.
A request the framework itself could not bind becomes \id{RQ1} per field under
\emph{request will not read}, naming each field \emph{as the wire spells it}:
\id{region\_id}, not \id{regionId}.

\section{Faults of the Request Itself}

Six rules are about the request rather than the map, and they exist because the
uniform shape held everywhere except at the door. Three are worth knowing by name.

\id{RQ1} is \emph{the document could not be read}: a 400, carrying the field where
the reader knew one.

\id{RQ3} is \emph{a field went unread}, and it is a complaint on a success rather
than a refusal. It is the rule that catches a typo: \id{roof.pich},
\id{wings[1].ptch}. The check walks down the posted value, matching by the name the
\emph{document} uses rather than the name the C\# type uses, and it runs over the
body as posted, before any merge. This is the studio's answer to the failure mode
where an agent posts a field the reader silently drops and believes it took effect.

\id{RQ5} is \emph{the request conflicts with what is stored}: a 409, and the one
refusal where nothing is wrong with the request at all. Every document carries a
revision, a read answers it as an \id{ETag}, and a write may state it back as an
\id{If-Match}. A request with no \id{If-Match} states no precondition and writes.
Protection is opted into by having read first, which is the only way it can mean
anything.

\section{One Question at Four Grains}

The most interesting structural property of the rule system is that the same
question recurs at four resolutions, with a different id at each, and each id can
only see what its own level knows.

Reachability is the worked example.

\begin{table}[htbp]
\caption{One question, \emph{can this be reached?}, asked at four grains. Read
down the column: a board can pass every gate above and still read a third dead.}
\label{tab:one-question}
\small
\begin{tabularx}{\textwidth}{l L{3.6cm} Y}
\toprule
\textbf{Rule} & \textbf{Grain} & \textbf{The question it asks} \\
\midrule
\id{WX6} & one piece of the plan & can a door be cut into this room at all \\
\id{PL9} & the whole plan & can a capturing team's spawn reach this wool \\
\id{EX1} & the built world & is the ground the rasterizer produced connected \\
\id{DR-PASS} & one building & does this building leave eight blocks to get past it \\
\id{DR-SLOPE} & one footprint & is this building's footprint level enough \\
(none) & a measurement & is reachable ground ground any journey actually uses \\
\bottomrule
\end{tabularx}
\end{table}

The documentation draws the conclusion, and it is the best argument in the system
for why there are four documents rather than one:

\begin{quote}
A finer grain exists only because a coarser one does \ldots Collapsing the four
documents into one would not simplify the rules. It would remove the vocabulary
each rule uses to say precisely what is wrong.
\end{quote}

Ids are grouped by what they are \emph{about}, never by which gate happens to ask.
The same objective rule is asked at the compile gate against a plan's pieces and
at the export gate against the ground the rasterizer produced, and a rule that
changed its name between the two would be two rules.

\section{Bands and Targets}

\Cref{ch:division} stated this as a finding; here is the mechanism. Rules divide
by \emph{kind}, and the two kinds most often confused are the target and the band.

A \textbf{target} is prescriptive and per-request: this compose wants exactly two
frontline runs, connected. A \textbf{band} is descriptive and measured off the
corpus: authored maps run one to seven frontline runs. The distinction is load
bearing, and the documentation states the failure precisely. A band pressed into
service as a target silently turns an observation about the corpus into a
requirement on the generator.

Several rules therefore read authored bands whose values depend on what the board
is played for. A frontline width cap of six to eight cells is the wool board's
figure; on a board played for cores or destroyables there is no width cap at all,
because a wide front feeds a wool lane and a destroy board has none. One run put
it as the most transferable thing it learned: three boards of experience on one
objective family taught it nothing about the next.

This is why \route{GET /api/rules/terms} exists. A driver asks what band is
enforced on \emph{this} board rather than writing down a number that was true of
the last one. The \ScoredTermCount{} scored terms and their bands are \cref{app:terms}.

\section{Nine Representative Rules}

The law is large and mostly technical, but a handful of its rules carry most of
what a reader needs to know about what the studio thinks a map is. These are
quoted as the law states them.

\paragraph{The corridor ladder.} \id{G2} fixes the whole board's scale:
the map's corridor width is the size band's, stated in blocks (12 at nano, 14 at
micro, 16 at milli and centi), with the wool approach one rung under it.
Measured over 359 built CTW maps.

\paragraph{The jump, not the bridge.} \id{WL12} is the newest wool rule and the
one an adapting agent breaks most often:

\begin{quote}
A bay or a hole beside a goal is at least 16 blocks across. Negative space is
crossed by \textbf{jumping} long before it is crossed by building \ldots Stated in
\textbf{blocks}, never in cells: a floor stated as a cell count moves with the
grid scale, and a jump does not care what the grid was.
\end{quote}

The floor is 16 blocks where the space touches a wool room or spawn and 12 where
it touches neither, because a hole in a team's own ground is crossed on purpose.

\paragraph{Holes are the rotation device.} \id{CT8} is the clearest statement in
the law of a rule that exists for a gameplay reason rather than a geometric one.
The closure encloses internal void pockets in every seed measured, and the
author's reason is that a loop around a hole gives alternative routes between
lanes without retreating through a single chokepoint. It is not mandatory, since
not every map has to have one though most do, so a holeless plan is flagged and
never blocked.

\paragraph{The strait.} \id{CT12}: on a two-team wool board the direct crossing
between the two team islands spans 15 to 40 blocks. Closer is two halves all but
merged; farther is a crossing nobody bridges.

\paragraph{The unwalkable seam.} \id{EL1}: a land seam of two blocks or more is
ground a player cannot walk up. The consequence is what the rule is for. On a
board built from the authored height palette every non-flush seam is a step nobody
walks, so the seam wants a ramp or a flight carved into it, and the check's output
is the work list for that pass rather than a fault to redraw around.

\paragraph{The mid never touches a wool.} \id{BZ6}, short and absolute: otherwise
players bridge freely through the mid straight into the point and there is no
gameplay direction.

\paragraph{A zone is bounded by the ground it docks.} \id{BZ9} forbids the
\emph{overhang}: a stretch carried beyond the last ground it meets, which connects
nothing and reads as no man's land a player walks into and stops.

\paragraph{The defence triangle.} \id{WL10} bans a specific failure: a front-near
wool whose spawn is far, beside a back wool whose spawn is adjacent, so that the
defender always guards the wrong door.

\paragraph{The destroy goal's ratio.} \id{GO1}: a destroy goal sits three to four
times as far from the enemy's spawn as from its own, measured by the walk over the
surface. Under the band the goal falls to a rush before a defence can form; over
it the attack is a march across the whole board into a set defence.

\begin{ruling}{Distance Is Measured as a Walk}
A distance between two things on a board is the walk between them over the
walkable surface. It is eight-connected, a diagonal costs what a player walks
rather than two steps, a climb is charged the blocks it takes to place, a fall is
counted but not charged, and the route goes around voids. Never the straight line through the air. Every
distance in every rule is this.
\end{ruling}
